% Teto: 400 palavras.
% Texto: seção 5 do .docx do PR 60, que ainda é só a marca do que escrever.
\section{Conclusão}
\label{sec:conclusao}

Este trabalho apresentou o SentryIoT, um sistema de NDR aberto para redes IoT que junta um classificador Random Forest a três agentes de LLM em um modelo local, e o avaliou como prova de conceito nos pcaps do CICIoT2023, na rede de uma casa com dispositivos IoT e em um teste com o Kali Linux. A detecção (P1) foi demonstrada nos ataques volumosos e só em parte nos demais: o classificador reconhece mais de 98\% das janelas de DDoS, DoS e Mirai, mas fica entre 34,4\% e 55,0\% em varredura, ataques Web e força bruta. Com a captura inteira fora do treino, o recall cai em todas as categorias com mais de um pcap; ainda assim, fora do treino, 39 das 40 capturas de ataque abriram incidente, com a regra de 4 janelas escolhida nesses mesmos testes. A resposta (P2) foi demonstrada no conjunto medido: a triagem reconheceu os 19 incidentes verdadeiros, e os agentes aplicaram uma ação aceitável no alvo certo em 16 deles e agiram em 1 dos 9 incidentes falsos, com uma limitação de taxa, de risco baixo, sobre um endereço que não era atacante, sem que nenhum endereço da rede chegasse ao agente de decisão nem que alguma ação atingisse o gateway ou os serviços críticos. Quanto ao custo e à rapidez (P3), nas execuções com ação aplicada, a remediação veio, na mediana, 60,86~s depois da primeira das janelas de ataque que abriram o incidente, e cada execução usou, na mediana, 64.206 tokens de entrada no modelo local, sem cobrança por chamada.

Esses resultados indicam que, nesta prova de conceito, um modelo local deu conta da triagem e da resposta de ataques ruidosos quando as garantias de segurança ficam fora do modelo, no servidor: na rodada final, a política de risco mandou para aprovação os bloqueios sobre endereços que não eram atacantes, e o agente de decisão trabalhou só com rótulos no lugar dos endereços; nos experimentos, quem aprovou foi uma regra que conhece os atacantes. Como cada análise, recomendação e ação fica gravada, quem opera a rede pode conferir por que o sistema agiu antes de aprovar e depois de executar.

A direção principal de trabalho futuro é medir o efeito real das ações em uma rede, com captura ao vivo, já que a prova de conceito simula a aplicação e a verificação. Outra é testar a resistência do próprio modelo a instruções maliciosas com um experimento de tool poisoning, porque os testes deste trabalho cobrem só o servidor. Na detecção, o recall de varredura, ataques Web e força bruta e a queda com capturas que não entraram no treino pedem dados de mais redes e mais capturas por técnica. Por fim, a arquitetura já permite trocar o modelo do agente de decisão por um modelo de fronteira, e comparar os dois em qualidade, custo e tempo mostraria quanto a resposta ganha com um modelo maior.
